\chapter{Regression Conclusions}
\label{chap:regression-conclusions}

\section{Main Findings}
\label{sec:reg-main-findings}

\subsection{Findings for Linear Regression}
Linear regression represents the most stable application for synthetic data generators. In continuous scenarios, parametric methods (Normal) generally attain somewhat higher CIO, while CART remains usable but tends to produce slightly wider confidence intervals. The number of predictors and correlation are the main sources of utility degradation.

\subsection{Findings for Logistic Regression}
Logistic regression is significantly more sensitive to synthesis artifacts. Small sample sizes ($N=100$) lead to low CIO and high bias, regardless of the generator. CART, in particular, often underestimates uncertainty (narrower CIs) in logistic outcomes, which is a noteworthy failure mode for inferential research.

\subsection{Findings that are Consistent Across Both Models}
\begin{itemize}
    \item \textbf{Sample Size:} Larger samples improve inferential fidelity, especially in the logistic setting.
    \item \textbf{Dimensionality:} Both models suffer as the number of predictors ($p$) increases.
    \item \textbf{Correlation:} High correlation ($\rho$) consistently degrades inferential fidelity.
    \item \textbf{Relative Generator Performance:} In scenarios where the data-generating process is compatible with the assumed parametric form, parametric synthesis (Normal, LogReg) generally performs somewhat better than CART for inferential preservation.
\end{itemize}

Across the simulated designs, inferential performance improved as sample size increased and deteriorated as the number of predictors or the predictor correlation increased. Linear regression remained comparatively stable, whereas logistic regression was markedly more fragile in small-sample and high-complexity settings. Taken together, these patterns indicate that generator performance depends not only on the synthesis method, but also on how demanding the inferential setting is.

\section{Comparison with Existing Literature}
Our findings align with and extend the work of researchers who have evaluated synthetic data in real-world settings (e.g., \cite{nowok2016, raab2016}). While those studies often highlight CART's flexibility for preserving complex univariate distributions, our systematic simulations suggest that for \textbf{inferential quantities} (regression coefficients), that flexibility can also be accompanied by structural artifacts that reduce CIO. This reinforces the idea that predictive utility does not always translate directly into inferential utility.

Specifically, compared to the ``baseline'' study using real data from the UK longitudinal studies, our simulations provide a clearer picture of possible \emph{mechanisms} of failure. The results suggest that ``synthetic bias'' is not merely a result of poor marginal matching, but can also emerge from the difficulty of capturing multi-dimensional conditional probabilities in sparse ($N/p$) regimes.

\section{Practical Recommendations}
\label{sec:reg-recommendations}

\subsection{Which Generators Perform Best Overall}
For researchers whose data aligns reasonably well with standard parametric assumptions (linear, normal, logistic), \textbf{Parametric Generators} are often a reasonable starting point for preserving inferential results.

\subsection{Which Generators Perform Best Under Specific Conditions}

Table~\ref{tab:generator_recommendations} provides a concise reference matrix for selecting a synthesis generator based on the data profile and the intended regression model.

\begin{table}[H]
  \centering
  \caption[Generator Recommendations by Scenario]{Synthesis generator recommendations matrix for inferential utility.}
  \label{tab:generator_recommendations}
  \renewcommand{\arraystretch}{1.3}
  \begin{tabular}{@{}lll@{}}
    \toprule
    \textbf{Data Condition} & \textbf{Linear Regression} & \textbf{Logistic Regression} \\ \midrule
    \textbf{Large Sample ($N \ge 1000$)} & Parametric (Normal) or CART & Parametric (LogReg) \\
    \textbf{Small Sample ($N \le 100$)} & Parametric (Normal) & Parametric (LogReg) \\
    \textbf{High Dimensionality ($p \ge 10$)} & Parametric (Normal) & \textit{High failure rate for all} \\
    \textbf{High Correlation ($\rho \ge 0.6$)} & Parametric (Normal) & Parametric (LogReg) \\ \bottomrule
  \end{tabular}
  \source{Compiled by the author}
\end{table}

\begin{itemize}
    \item \textbf{High N, Low p:} CART and Parametric are both viable.
    \item \textbf{Small N:} Parametric synthesis is generally preferable, as CART becomes less stable.
    \item \textbf{Binary Outcomes:} LogReg synthesis is typically preferable; CART should be used cautiously when inference on the logistic model is the main goal.
\end{itemize}

\subsection{Guidance for Researchers Choosing a Generator}
Researchers should prioritize generators that match their intended analytical model. If a linear regression is planned for the final analysis, a linear (Normal) synthesis arm is often an appropriate baseline to evaluate first.

\section{Limitations}
\label{sec:reg-limitations}
\begin{itemize}
    \item \textbf{Simulation Focus:} All datasets were generated from known distributions. Real-world non-linearities or ``messy'' data may change the relative performance of CART.
    \item \textbf{Metric Focus:} We focused on CIO. Other metrics (e.g., coverage, MSE) might reveal different tradeoffs.
\end{itemize}

\section{Future Work}
\label{sec:reg-future-work}
Future research should explore the intersection of \textbf{Differential Privacy} and inferential utility \cite{dwork2006, jordon2018}. Understanding how the injection of noise (e.g., $\epsilon$-privacy) further degrades CIO across different generators is a critical next step for the field. Expanding this research into settings combining differential privacy guarantees with rigorous disclosure analysis \cite{abowd2003} and comprehensive synthetic dataset evaluation \cite{bellovin2019, bowen2020, hsu2021} will provide practitioners with more robust frameworks. Finally, comparing these synthesis approaches against traditional multiple imputation techniques \cite{rubin1987, raghunathan2003} could yield further insights into the trade-offs between privacy and statistical validity.
